\documentclass[11pt,a4paper]{article}

\usepackage[utf8]{inputenc}
\usepackage[T1]{fontenc}
\usepackage{amsmath,amssymb}
\usepackage{booktabs}
\usepackage{longtable}
\usepackage{array}
\usepackage{hyperref}
\usepackage{geometry}
\usepackage{enumitem}
\usepackage{listings}
\usepackage{xcolor}

\geometry{margin=1in}
\hypersetup{colorlinks=true, linkcolor=blue, urlcolor=blue, citecolor=blue}

\lstset{
  basicstyle=\ttfamily\small,
  breaklines=true,
  frame=single,
  backgroundcolor=\color{gray!8},
  columns=fullflexible
}

\newcommand{\COtwo}{CO\textsubscript{2}}

\title{Calibration of Shortfall, Wastage, and Transmission\\
Operating Costs in the Nested WECC--California Model}
\author{ASTR2026 transmission layer}
\date{September 2026}

\begin{document}
\maketitle

\section{Purpose}

The nested WECC--California linear program prices three quantities that are not
physical generation: \emph{shortfall} (unserved demand), \emph{wastage}
(curtailed surplus), and \emph{transmission operating cost} (per-corridor flow
cost). These are modelling parameters rather than measured data, and their
values determine how the optimiser trades unserved energy, curtailment, and
wheeling against real generation cost.

This document records the benchmarks used to set them, the derivations, and the
diagnosis that motivated the recalibration.

\section{Unit convention}

All costs in the model are expressed in dollars per joule, \$/J. Published
benchmarks are almost universally quoted in dollars per megawatt hour. The
conversion is
\begin{equation}
1~\mathrm{MWh} = 10^{6}~\mathrm{W} \times 3600~\mathrm{s}
              = 3.6\times10^{9}~\mathrm{J},
\qquad
\left[\frac{\$}{\mathrm{MWh}}\right]
  = \left[\frac{\$}{\mathrm{J}}\right] \times 3.6\times 10^{9}.
\end{equation}
Failing to apply this factor is what produced the miscalibration described in
Section~\ref{sec:diagnosis}.

\section{Diagnosis: the miscalibration}
\label{sec:diagnosis}

The generator operating costs carried in the WECC graph are correctly scaled:
nuclear at $3.09\times10^{-9}$~\$/J is \$11.13/MWh, and the most expensive unit
in the fleet at $1.07\times10^{-8}$~\$/J is \$38.52/MWh. Both are realistic.
The penalty parameters were not on the same scale.

\begin{table}[h]
\centering
\caption{Pre-calibration parameter values expressed in \$/MWh.}
\begin{tabular}{lrrl}
\toprule
Parameter & \$/J (previous) & = \$/MWh & Benchmark \\
\midrule
\texttt{shortfall\_cost} & $1\times10^{-2}$ & 36{,}000{,}000 & \$10{,}000/MWh \\
\texttt{wastage\_cost}   & $1\times10^{-4}$ & 360{,}000      & \$0--25/MWh \\
Transmission (meso)      & $0$              & $0$            & $>0$ required \\
\bottomrule
\end{tabular}
\end{table}

Two consequences were measured on the 672-hour S0 solve.

\paragraph{The objective was almost entirely penalty.}
With shortfall of 36.523~GWh and wastage of 592.954~GWh:
\begin{align}
\text{shortfall penalty} &= 36.523~\text{GWh} \times \$3.6\times10^{7}/\text{MWh} = \$1.315\times10^{12}, \\
\text{wastage penalty}   &= 592.954~\text{GWh} \times \$3.6\times10^{5}/\text{MWh} = \$0.213\times10^{12}, \\
\text{sum}               &= \$1.528\times10^{12},
\end{align}
against a total objective of $\$1.528569\times10^{12}$. More than 99.9\% of the
objective was these two penalties. Real generation cost --- which determines
which plants run, and therefore \COtwo{} --- accounted for roughly 0.04\%.

\paragraph{The LP was badly conditioned.}
Objective coefficients spanned $1.98\times10^{-13}$ to $1\times10^{-2}$, about
eleven orders of magnitude. Gurobi reported coefficient-range warnings on every
solve, and the S3 scenario terminated with \texttt{Numerical trouble
encountered}.

Because generation cost was numerically negligible, the dispatch decision was
close to degenerate: the solver had little incentive to distinguish between
generators. This is consistent with the observed sensitivity of per-asset
results to the Gurobi \texttt{Crossover} setting, in which barrier-only and
crossover solutions agreed on the objective to 0.04\% but differed by a factor
of 2.6 in \COtwo{} and 3.5 in fossil generation.

\section{Calibrated values}

\begin{longtable}{p{2.4cm}rrp{4.4cm}p{2.6cm}}
\caption{Calibrated parameters, benchmarks, and sources.}\\
\toprule
Parameter & Value (\$/J) & = \$/MWh & Benchmark and rationale & Source \\
\midrule
\endfirsthead
\toprule
Parameter & Value (\$/J) & = \$/MWh & Benchmark and rationale & Source \\
\midrule
\endhead
\bottomrule
\endfoot

\texttt{shortfall\_cost} (VOLL)
& $2.78\times10^{-6}$ & 10{,}000
& Value of lost load. \$10{,}000/MWh is the default VOLL in the GenX capacity
  expansion model and sits in the lower-middle of the \$9{,}000--\$45{,}000/MWh
  range reported for developed economies. ERCOT adopted \$35{,}000/MWh; EU
  member state values range from EUR 4{,}000 to EUR 69{,}000/MWh. Retained at
  \$10{,}000/MWh as a conservative, widely used default; still about
  $260\times$ the priciest generator, so load is served whenever physically
  possible.
& \cite{genx,brattle2024,ercot2013,miso2022,euets} \\
\addlinespace

\texttt{wastage\_cost} (curtailment)
& $2.78\times10^{-10}$ & 1.00
& Curtailment of surplus is conventionally valued at \$0/MWh. A strictly
  positive value is nonetheless required here as a tie-breaker: at exactly
  zero, corridor arc pairs form a flat direction in the objective that an
  interior-point method smears across. \$1/MWh is the low end of the
  curtailment range and far below any generator's cost, so it breaks ties
  without steering dispatch.
& \cite{flexpower,rff} \\
\addlinespace

Transmission operating cost (meso corridors)
& $1.0\times10^{-10}$ & 0.36
& Published wheeling charges are about \$2/MWh, but price a whole
  point-to-point transaction rather than each line crossed. The BA tier charges
  transaction-scale rates (\$8/MWh on 40 of 68 lines) because BA paths are few
  hops; meso paths cross many of about 5{,}200 substation corridors, so a
  transaction-scale per-hop rate would dominate merit order. \$0.36/MWh per
  corridor puts a ten-hop intra-CA path near \$3.6/MWh --- the right order of
  magnitude for a transaction, and well below gas (about \$38/MWh).
& \cite{brown2024,lbnl2019} \\
\end{longtable}

\subsection{Why transmission cost must be strictly positive}

Each corridor is represented as two opposing directed arcs, because the
\texttt{Transmission} class holds a single non-negative flow variable and
applies its efficiency only on the receiving side. No linear constraint can
forbid both arcs of a pair from being active simultaneously. With a zero flow
cost, a corridor pair is therefore a perfectly flat direction in the objective:
energy may be cycled $A \rightarrow B \rightarrow A$ indefinitely, burning
$(1-\eta)$ per pass at no cost.

A strictly positive per-unit flow cost removes the flatness. Cancelling the
common component $\delta$ of any bidirectional pair strictly lowers cost by
$2\delta c$ \emph{and} returns the round-trip losses $(1-\eta)\delta$ to both
end nodes, so simultaneous bidirectional flow can never be optimal. The base
WECC graph achieves the same effect by being lossless ($\eta = 1$) while
charging $2.222\times10^{-9}$~\$/J on most BA lines.

Introducing the cost reduced gross transmission flow in the crossover solution
from 293.3~TWh to 28.7~TWh, and corridors carrying flow in both directions
simultaneously from 1{,}097 to 268, while changing \COtwo{} by 0.1\%.

\section{Expected effects of recalibration}

\begin{enumerate}[leftmargin=*]
  \item The objective coefficient range narrows from about eleven to about
        seven orders of magnitude, improving conditioning.
  \item Generation cost becomes the dominant term, so dispatch is driven by
        generator economics rather than by penalty parameters. This is expected
        to reduce, though not necessarily eliminate, the degeneracy underlying
        the \texttt{Crossover} sensitivity.
  \item Objective values become interpretable in dollars.
  \item Shortfall at \$10{,}000/MWh remains about $260\times$ the most
        expensive generator, so unserved energy appears only where demand is
        genuinely infeasible to serve, not as an economic substitute for
        generation.
\end{enumerate}

\section{Recommended sensitivity runs}

\begin{itemize}[leftmargin=*]
  \item \textbf{VOLL:} \$35{,}000/MWh ($9.72\times10^{-6}$~\$/J), the ERCOT
        value, to test whether conclusions depend on the penalty for unserved
        energy.
  \item \textbf{Curtailment:} \$25/MWh ($6.94\times10^{-9}$~\$/J), the wind
        production tax credit opportunity cost, which is the economically
        grounded alternative to a pure tie-breaker.
  \item \textbf{Transmission:} \$1/MWh per corridor
        ($2.78\times10^{-10}$~\$/J) to test sensitivity of flow patterns to the
        wheeling assumption.
\end{itemize}

\section{Where these are set}

\begin{lstlisting}
ev_charging_project/config.py
    NETWORK_KW['shortfall_cost'] = 10_000 / 3.6e9   # 2.78e-6  $/J
    NETWORK_KW['wastage_cost']   =      1 / 3.6e9   # 2.78e-10 $/J

Distribution-Grid-EV-CA/09_02_nest_meso_in_good.py
    TRANSMISSION_OPERATING_COST  =   0.36 / 3.6e9   # 1.0e-10  $/J
\end{lstlisting}

\texttt{10\_01\_run\_scenarios\_S0\_S3.py} no longer overrides the shortfall and
wastage costs; \texttt{config.py} is the single source of truth.

\begin{thebibliography}{9}

\bibitem{genx}
GenX Project.
\emph{GenX: a configurable power system capacity expansion model for studying
low-carbon energy futures.} MIT Energy Initiative.
\url{https://github.com/GenXProject/GenX.jl};
\url{https://genx.mit.edu}.
Default value of lost load \$10{,}000/MWh.

\bibitem{brattle2024}
The Brattle Group.
\emph{Value of Lost Load Study for the ERCOT Region}, Project No.\ 55837, 2024.
\url{https://www.brattle.com/wp-content/uploads/2024/09/Value-of-Lost-Load-Study-for-the-ERCOT-Region.pdf}

\bibitem{ercot2013}
ERCOT.
\emph{Estimating the Value of Lost Load: Literature Review and Macroeconomic
Analysis}, 2013.
\url{https://www.ercot.com/files/docs/2013/06/19/ercot_valueoflostload_literaturereviewandmacroeconomic.pdf}

\bibitem{miso2022}
MISO.
\emph{Value of Lost Load (VOLL) Development and Use}, RECBWG, August 2022.
\url{https://cdn.misoenergy.org/}

\bibitem{euets}
Emissions-EUETS.
\emph{Value of Lost Load (VoLL)}, internal electricity market glossary.
\url{https://www.emissions-euets.com/internal-electricity-market-glossary/966-value-of-lost-load-voll}
EU member state VOLL values.

\bibitem{flexpower}
Flex Power.
\emph{The Economics of Curtailing Renewables like Wind and Solar}.
\url{https://flex-power.energy/energyblog/the-economics-of-curtailing-renewables/}
Zero-cost curtailment convention and PTC-driven negative bidding.

\bibitem{rff}
Resources for the Future.
\emph{What Are the Costs and Values of Wind and Solar Power? How Are They
Changing?}
\url{https://www.resources.org/common-resources/what-are-costs-and-values-wind-and-solar-power-how-are-they-changing/}
Production tax credit of about \$25/MWh as curtailment opportunity cost.

\bibitem{brown2024}
Brown, P.\ et al.
\emph{Optimal transmission expansion modestly reduces decarbonization costs of
U.S. electricity}, arXiv:2402.14189.
\url{https://arxiv.org/pdf/2402.14189}
Uses a \$2.00/MWh wheeling rate.

\bibitem{lbnl2019}
Lawrence Berkeley National Laboratory.
\emph{Improving estimates of transmission capital costs for utility-scale wind
and solar projects to inform renewable energy policy}, 2019.
\url{https://eta-publications.lbl.gov/sites/default/files/td_costs_formatted_final.pdf}

\end{thebibliography}

\end{document}
